% Deck 1 of 4: Parts I and II of AI for Biodiversity: A Primer.
% Build: latexmk -pdf main.tex   (pdflatex + biber)
\documentclass[11pt,aspectratio=169]{beamer}
% ---------------------------------------------------------------------------
% preamble.tex --- shared Beamer style for the four workshop decks
%
% Loaded by each deck with % ---------------------------------------------------------------------------
% preamble.tex --- shared Beamer style for the four workshop decks
%
% Loaded by each deck with % ---------------------------------------------------------------------------
% preamble.tex --- shared Beamer style for the four workshop decks
%
% Loaded by each deck with \input{../common/preamble}. Same fonts, colours
% and bibliography settings as the book (book/preamble.tex): Latin Modern,
% black on white, one grey for code, one accent colour for what is untrusted
% or wrong. Built with pdflatex + biber, like the book.
% ---------------------------------------------------------------------------

\usepackage[T1]{fontenc}
\usepackage[utf8]{inputenc}
\usepackage{lmodern}
\usepackage{microtype}
\usepackage{booktabs}
\usepackage{tabularx}
\usepackage{listings}
\usepackage{xcolor}
\usepackage{csquotes}
\usepackage{graphicx}

% Figures: the book's figures as the 300 dpi PNGs the book itself includes
% (book/figures/png, tracked in git and current; figures/pdf is not rebuilt
% after every change), and the first pages of the papers
% (common/screenshots). Both are found by name without an extension.
\graphicspath{{../../book/figures/png/}{../common/screenshots/}}
\DeclareGraphicsExtensions{.png,.pdf}

% --- look ------------------------------------------------------------------
\usetheme{default}
\usefonttheme{serif}
\definecolor{accent}{HTML}{9B2D20}
\definecolor{codebg}{gray}{0.96}
\definecolor{rule}{gray}{0.75}
\definecolor{dim}{gray}{0.45}
\setbeamercolor{structure}{fg=black}
\setbeamercolor{normal text}{fg=black,bg=white}
\setbeamercolor{alerted text}{fg=accent}
\setbeamercolor{section in toc}{fg=black}
\setbeamercolor{footline}{fg=dim}
\setbeamerfont{title}{series=\bfseries,size=\Large}
\setbeamerfont{subtitle}{series=\mdseries,size=\normalsize}
\setbeamerfont{frametitle}{series=\bfseries,size=\large}
\setbeamerfont{footline}{size=\scriptsize}
\setbeamertemplate{navigation symbols}{}
\setbeamertemplate{itemize items}[circle]
\setbeamertemplate{enumerate items}[default]
\setbeamertemplate{caption}{\scriptsize\insertcaption}
\setbeamertemplate{frametitle}{%
  \vskip0.6em\usebeamerfont{frametitle}\insertframetitle\par
  \vskip0.2em{\color{rule}\hrule height 0.4pt}\vskip0.4em}
\setbeamertemplate{footline}{%
  \hbox to \paperwidth{\hspace{1.2em}\usebeamerfont{footline}\usebeamercolor[fg]{footline}%
    \insertshorttitle\ --- \insertsection\hfill
    \insertframenumber/\inserttotalframenumber\hspace{1.2em}}\vskip0.6em}
\setbeamersize{text margin left=1.4em,text margin right=1.4em}
% An outline frame at the start of each section, the current one in black.
\AtBeginSection[]{\begin{frame}{Outline}\tableofcontents[currentsection]\end{frame}}
% A frame that breaks over pages is numbered only from its second page.
\setbeamertemplate{frametitle continuation}[from second]

% --- code listings (as in the book) ----------------------------------------
\lstset{
  basicstyle=\ttfamily\footnotesize,
  backgroundcolor=\color{codebg},
  frame=none,
  breaklines=true,
  breakatwhitespace=true,
  columns=fullflexible,
  keepspaces=true,
  upquote=true,
  showstringspaces=false,
  xleftmargin=0.6em,
  framexleftmargin=0.6em,
  aboveskip=0.5em,
  belowskip=0.5em,
  literate={~}{{\textasciitilde}}1
           {>}{{\textgreater}}1
           {<}{{\textless}}1
           {ö}{{\"o}}1 {ü}{{\"u}}1 {é}{{\'e}}1,
}
\lstdefinestyle{shell}{language=bash}

% --- bibliography (same settings as the book) ------------------------------
\usepackage[
  backend=biber,
  style=authoryear,
  sorting=nyt,
  maxbibnames=6,
  minbibnames=3,
  maxcitenames=2,
  giveninits=true,
  uniquename=init,
  doi=true,
  url=true,
  isbn=false,
  eprint=true,
]{biblatex}
\addbibresource{../../book/references.bib}
\renewcommand*{\bibfont}{\scriptsize}
\setlength{\bibitemsep}{0.25em}
\setbeamertemplate{bibliography item}{}   % no icon before each entry

% --- shorthand -------------------------------------------------------------
\newcommand{\code}[1]{\texttt{#1}}
\newcommand{\term}[1]{\emph{#1}}
% A key point, as in the book: no box.
\newcommand{\kp}[1]{\par\medskip{\bfseries Key point.}\enspace\emph{#1}\par}
% \fig[height]{name}: one of the book's figures, scaled to fit the frame.
\newcommand{\fig}[2][0.72\textheight]{%
  \centering\includegraphics[height=#1,width=\textwidth,keepaspectratio]{#2}\par}
% \screenshot{key}: the first page of a paper, with a thin frame.
\newcommand{\screenshot}[2][0.66\textheight]{%
  \setlength{\fboxsep}{0pt}\setlength{\fboxrule}{0.3pt}%
  \fbox{\includegraphics[height=#1,width=\linewidth,keepaspectratio]{#2}}}
% \paperslide{key}{text}: screenshot on the left, what the book takes from the
% paper on the right. Use inside a frame; add \source after it.
\newcommand{\paperslide}[2]{%
  \begin{columns}[T]
    \begin{column}{0.40\textwidth}\centering\screenshot{#1}\end{column}
    \begin{column}{0.58\textwidth}\small #2\end{column}
  \end{columns}}
% Notes in references.bib record how a web source was saved and dated; they
% belong in the reference list, not on a \source line, so cite commands
% (including \fullcite) leave them out.
\AtEveryCitekey{\clearfield{note}}
% \source{...}: a small grey line at the foot of a frame, for the source.
\newcommand{\source}[1]{\vskip0pt plus 1filll\par{\scriptsize\color{dim}#1\par}}
. Same fonts, colours
% and bibliography settings as the book (book/preamble.tex): Latin Modern,
% black on white, one grey for code, one accent colour for what is untrusted
% or wrong. Built with pdflatex + biber, like the book.
% ---------------------------------------------------------------------------

\usepackage[T1]{fontenc}
\usepackage[utf8]{inputenc}
\usepackage{lmodern}
\usepackage{microtype}
\usepackage{booktabs}
\usepackage{tabularx}
\usepackage{listings}
\usepackage{xcolor}
\usepackage{csquotes}
\usepackage{graphicx}

% Figures: the book's figures as the 300 dpi PNGs the book itself includes
% (book/figures/png, tracked in git and current; figures/pdf is not rebuilt
% after every change), and the first pages of the papers
% (common/screenshots). Both are found by name without an extension.
\graphicspath{{../../book/figures/png/}{../common/screenshots/}}
\DeclareGraphicsExtensions{.png,.pdf}

% --- look ------------------------------------------------------------------
\usetheme{default}
\usefonttheme{serif}
\definecolor{accent}{HTML}{9B2D20}
\definecolor{codebg}{gray}{0.96}
\definecolor{rule}{gray}{0.75}
\definecolor{dim}{gray}{0.45}
\setbeamercolor{structure}{fg=black}
\setbeamercolor{normal text}{fg=black,bg=white}
\setbeamercolor{alerted text}{fg=accent}
\setbeamercolor{section in toc}{fg=black}
\setbeamercolor{footline}{fg=dim}
\setbeamerfont{title}{series=\bfseries,size=\Large}
\setbeamerfont{subtitle}{series=\mdseries,size=\normalsize}
\setbeamerfont{frametitle}{series=\bfseries,size=\large}
\setbeamerfont{footline}{size=\scriptsize}
\setbeamertemplate{navigation symbols}{}
\setbeamertemplate{itemize items}[circle]
\setbeamertemplate{enumerate items}[default]
\setbeamertemplate{caption}{\scriptsize\insertcaption}
\setbeamertemplate{frametitle}{%
  \vskip0.6em\usebeamerfont{frametitle}\insertframetitle\par
  \vskip0.2em{\color{rule}\hrule height 0.4pt}\vskip0.4em}
\setbeamertemplate{footline}{%
  \hbox to \paperwidth{\hspace{1.2em}\usebeamerfont{footline}\usebeamercolor[fg]{footline}%
    \insertshorttitle\ --- \insertsection\hfill
    \insertframenumber/\inserttotalframenumber\hspace{1.2em}}\vskip0.6em}
\setbeamersize{text margin left=1.4em,text margin right=1.4em}
% An outline frame at the start of each section, the current one in black.
\AtBeginSection[]{\begin{frame}{Outline}\tableofcontents[currentsection]\end{frame}}
% A frame that breaks over pages is numbered only from its second page.
\setbeamertemplate{frametitle continuation}[from second]

% --- code listings (as in the book) ----------------------------------------
\lstset{
  basicstyle=\ttfamily\footnotesize,
  backgroundcolor=\color{codebg},
  frame=none,
  breaklines=true,
  breakatwhitespace=true,
  columns=fullflexible,
  keepspaces=true,
  upquote=true,
  showstringspaces=false,
  xleftmargin=0.6em,
  framexleftmargin=0.6em,
  aboveskip=0.5em,
  belowskip=0.5em,
  literate={~}{{\textasciitilde}}1
           {>}{{\textgreater}}1
           {<}{{\textless}}1
           {ö}{{\"o}}1 {ü}{{\"u}}1 {é}{{\'e}}1,
}
\lstdefinestyle{shell}{language=bash}

% --- bibliography (same settings as the book) ------------------------------
\usepackage[
  backend=biber,
  style=authoryear,
  sorting=nyt,
  maxbibnames=6,
  minbibnames=3,
  maxcitenames=2,
  giveninits=true,
  uniquename=init,
  doi=true,
  url=true,
  isbn=false,
  eprint=true,
]{biblatex}
\addbibresource{../../book/references.bib}
\renewcommand*{\bibfont}{\scriptsize}
\setlength{\bibitemsep}{0.25em}
\setbeamertemplate{bibliography item}{}   % no icon before each entry

% --- shorthand -------------------------------------------------------------
\newcommand{\code}[1]{\texttt{#1}}
\newcommand{\term}[1]{\emph{#1}}
% A key point, as in the book: no box.
\newcommand{\kp}[1]{\par\medskip{\bfseries Key point.}\enspace\emph{#1}\par}
% \fig[height]{name}: one of the book's figures, scaled to fit the frame.
\newcommand{\fig}[2][0.72\textheight]{%
  \centering\includegraphics[height=#1,width=\textwidth,keepaspectratio]{#2}\par}
% \screenshot{key}: the first page of a paper, with a thin frame.
\newcommand{\screenshot}[2][0.66\textheight]{%
  \setlength{\fboxsep}{0pt}\setlength{\fboxrule}{0.3pt}%
  \fbox{\includegraphics[height=#1,width=\linewidth,keepaspectratio]{#2}}}
% \paperslide{key}{text}: screenshot on the left, what the book takes from the
% paper on the right. Use inside a frame; add \source after it.
\newcommand{\paperslide}[2]{%
  \begin{columns}[T]
    \begin{column}{0.40\textwidth}\centering\screenshot{#1}\end{column}
    \begin{column}{0.58\textwidth}\small #2\end{column}
  \end{columns}}
% Notes in references.bib record how a web source was saved and dated; they
% belong in the reference list, not on a \source line, so cite commands
% (including \fullcite) leave them out.
\AtEveryCitekey{\clearfield{note}}
% \source{...}: a small grey line at the foot of a frame, for the source.
\newcommand{\source}[1]{\vskip0pt plus 1filll\par{\scriptsize\color{dim}#1\par}}
. Same fonts, colours
% and bibliography settings as the book (book/preamble.tex): Latin Modern,
% black on white, one grey for code, one accent colour for what is untrusted
% or wrong. Built with pdflatex + biber, like the book.
% ---------------------------------------------------------------------------

\usepackage[T1]{fontenc}
\usepackage[utf8]{inputenc}
\usepackage{lmodern}
\usepackage{microtype}
\usepackage{booktabs}
\usepackage{tabularx}
\usepackage{listings}
\usepackage{xcolor}
\usepackage{csquotes}
\usepackage{graphicx}

% Figures: the book's figures as the 300 dpi PNGs the book itself includes
% (book/figures/png, tracked in git and current; figures/pdf is not rebuilt
% after every change), and the first pages of the papers
% (common/screenshots). Both are found by name without an extension.
\graphicspath{{../../book/figures/png/}{../common/screenshots/}}
\DeclareGraphicsExtensions{.png,.pdf}

% --- look ------------------------------------------------------------------
\usetheme{default}
\usefonttheme{serif}
\definecolor{accent}{HTML}{9B2D20}
\definecolor{codebg}{gray}{0.96}
\definecolor{rule}{gray}{0.75}
\definecolor{dim}{gray}{0.45}
\setbeamercolor{structure}{fg=black}
\setbeamercolor{normal text}{fg=black,bg=white}
\setbeamercolor{alerted text}{fg=accent}
\setbeamercolor{section in toc}{fg=black}
\setbeamercolor{footline}{fg=dim}
\setbeamerfont{title}{series=\bfseries,size=\Large}
\setbeamerfont{subtitle}{series=\mdseries,size=\normalsize}
\setbeamerfont{frametitle}{series=\bfseries,size=\large}
\setbeamerfont{footline}{size=\scriptsize}
\setbeamertemplate{navigation symbols}{}
\setbeamertemplate{itemize items}[circle]
\setbeamertemplate{enumerate items}[default]
\setbeamertemplate{caption}{\scriptsize\insertcaption}
\setbeamertemplate{frametitle}{%
  \vskip0.6em\usebeamerfont{frametitle}\insertframetitle\par
  \vskip0.2em{\color{rule}\hrule height 0.4pt}\vskip0.4em}
\setbeamertemplate{footline}{%
  \hbox to \paperwidth{\hspace{1.2em}\usebeamerfont{footline}\usebeamercolor[fg]{footline}%
    \insertshorttitle\ --- \insertsection\hfill
    \insertframenumber/\inserttotalframenumber\hspace{1.2em}}\vskip0.6em}
\setbeamersize{text margin left=1.4em,text margin right=1.4em}
% An outline frame at the start of each section, the current one in black.
\AtBeginSection[]{\begin{frame}{Outline}\tableofcontents[currentsection]\end{frame}}
% A frame that breaks over pages is numbered only from its second page.
\setbeamertemplate{frametitle continuation}[from second]

% --- code listings (as in the book) ----------------------------------------
\lstset{
  basicstyle=\ttfamily\footnotesize,
  backgroundcolor=\color{codebg},
  frame=none,
  breaklines=true,
  breakatwhitespace=true,
  columns=fullflexible,
  keepspaces=true,
  upquote=true,
  showstringspaces=false,
  xleftmargin=0.6em,
  framexleftmargin=0.6em,
  aboveskip=0.5em,
  belowskip=0.5em,
  literate={~}{{\textasciitilde}}1
           {>}{{\textgreater}}1
           {<}{{\textless}}1
           {ö}{{\"o}}1 {ü}{{\"u}}1 {é}{{\'e}}1,
}
\lstdefinestyle{shell}{language=bash}

% --- bibliography (same settings as the book) ------------------------------
\usepackage[
  backend=biber,
  style=authoryear,
  sorting=nyt,
  maxbibnames=6,
  minbibnames=3,
  maxcitenames=2,
  giveninits=true,
  uniquename=init,
  doi=true,
  url=true,
  isbn=false,
  eprint=true,
]{biblatex}
\addbibresource{../../book/references.bib}
\renewcommand*{\bibfont}{\scriptsize}
\setlength{\bibitemsep}{0.25em}
\setbeamertemplate{bibliography item}{}   % no icon before each entry

% --- shorthand -------------------------------------------------------------
\newcommand{\code}[1]{\texttt{#1}}
\newcommand{\term}[1]{\emph{#1}}
% A key point, as in the book: no box.
\newcommand{\kp}[1]{\par\medskip{\bfseries Key point.}\enspace\emph{#1}\par}
% \fig[height]{name}: one of the book's figures, scaled to fit the frame.
\newcommand{\fig}[2][0.72\textheight]{%
  \centering\includegraphics[height=#1,width=\textwidth,keepaspectratio]{#2}\par}
% \screenshot{key}: the first page of a paper, with a thin frame.
\newcommand{\screenshot}[2][0.66\textheight]{%
  \setlength{\fboxsep}{0pt}\setlength{\fboxrule}{0.3pt}%
  \fbox{\includegraphics[height=#1,width=\linewidth,keepaspectratio]{#2}}}
% \paperslide{key}{text}: screenshot on the left, what the book takes from the
% paper on the right. Use inside a frame; add \source after it.
\newcommand{\paperslide}[2]{%
  \begin{columns}[T]
    \begin{column}{0.40\textwidth}\centering\screenshot{#1}\end{column}
    \begin{column}{0.58\textwidth}\small #2\end{column}
  \end{columns}}
% Notes in references.bib record how a web source was saved and dated; they
% belong in the reference list, not on a \source line, so cite commands
% (including \fullcite) leave them out.
\AtEveryCitekey{\clearfield{note}}
% \source{...}: a small grey line at the foot of a frame, for the source.
\newcommand{\source}[1]{\vskip0pt plus 1filll\par{\scriptsize\color{dim}#1\par}}


\title[AI in Practice 1]{AI in Practice, or: How I Learned to Stop Worrying and Love the Agent}
\subtitle{A toolbox and worked examples for biodiversity researchers, from the book \emph{AI for Biodiversity: A Primer}\\[0.6em] 1. The landscape and the data (Chapters 1--6)}
\author{Savvas Paragkamian}
\institute{Hellenic Centre for Marine Research \\ KYBELE and ELIXIR-GR}
\date{HCMR workshop, November 2026}
\hypersetup{pdftitle={AI for Biodiversity: A Primer, deck 1},pdfauthor={Savvas Paragkamian}}

\begin{document}

\begin{frame}[plain]
  \titlepage
\end{frame}

\begin{frame}{Goals of this session}
  \begin{enumerate}
    \item Tell a model, a chatbot and an agent apart, and know why these
      systems started working this well only recently.
    \item Know what language models do well, how they fail, and what the
      ecology literature shows so far.
    \item Know what ``AI-ready'' hides, what makes biodiversity records fit for
      a model, and how to organise and protect your own data before an agent
      touches it.
  \end{enumerate}
  \kp{The machine takes the drudgery, and the judgement stays with the
    ecologist. That is the frame of the book.}
\end{frame}

% ===========================================================================
\section{Chatbots, models and agents, and why now (Chapters 1 and 2)}

\begin{frame}{How the field worked before computers}
  \small
  \begin{itemize}
    \item The founders worked on paper. The calculation each is remembered for
      took an evening, a weekend or a season; the record-keeping, arithmetic,
      reading and writing around it took years, and much of that was done by
      assistants.
    \item Chapter 1 opens with six portraits, each built from the four jobs
      this book is about: gathering data, calculating, finding what has been
      written, and writing it up.
    \item D'Arcy Thompson (1860--1948), R.\,A. Fisher (1890--1962), Claude
      Shannon (1916--2001), John Tukey (1915--2000), Robert May (1936--2020)
      and William Michener, who built the record-keeping the others did
      without.
    \item The tools were shared: paper slips and India-ink labels, printed
      tables, slide rules, the Millionaire calculator, human computers,
      punched cards, reprint postcards and typescripts; then GENSTAT, S and R,
      metadata and repositories.
  \end{itemize}
  \medskip
  Nobody in that story timed the clerical steps between decisions. The
  earliest timing the book has found is Licklider's, on the next slide.
\end{frame}

\begin{frame}{D'Arcy Thompson: fishery statistics and a fur-seal census}
  \paperslide{thompson1905aberdeen}{\footnotesize
    \begin{itemize}
      \item Scientific member of the Fishery Board for Scotland for over forty
        years; by his daughter's account, editor of the ICES statistical
        bulletin, 1902--1947 \parencite{lt1958scholarnaturalist,thompsonr1976dsb}.
      \item His 1905 tables of Aberdeen trawl landings came from the Board's
        reports; he noted what they could not show, such as time spent
        fishing \parencite{thompson1905aberdeen}.
      \item His one field census, the Bering Sea fur seals of 1896, was made
        alongside the American commission; he counted harems himself
        \parencite{thompson1897beringsea}.
      \item Calculation he left to others: ``I pretend to no mathematical
        skill''; Mackinnon and Ogilvie drew most of the diagrams of
        \emph{On Growth and Form} \parencite{thompson1917growthform}.
    \end{itemize}}
  % Short source line: \fullcite adds the URL and runs to three lines here.
  \source{\citeauthor{thompson1905aberdeen} (\citeyear{thompson1905aberdeen}).
    \citetitle{thompson1905aberdeen}.}
\end{frame}

\begin{frame}{R.\,A. Fisher at Rothamsted: ledgers and a Millionaire}
  \paperslide{fisher1926arrangement}{\footnotesize
    \begin{itemize}
      \item Hired in 1919 to re-examine seventy years of hand-kept
        field-experiment records, a task he called ``raking over the muck
        heap'' \parencite{yates1963fishermemoir}.
      \item A calculating machine was rented for him in 1919; in 1921 or 1922
        he bought a motor-driven Millionaire, its price said to exceed
        \pounds200 \parencite{gower2015rothamstedcomputing}.
      \item The routine arithmetic fell to one or two assistants and a handful
        of human computers, mostly women; a complex design took weeks or
        months \parencite{parolini2015rothamsted}.
      \item \emph{Statistical Methods for Research Workers} gave tables for
        fixed probabilities, so the reader never computed one
        \parencite{parolini2015rothamsted}; in 1926 the laboratory offered
        randomised layouts to any experimenter
        \parencite{fisher1926arrangement}.
    \end{itemize}}
  % Short source line: \fullcite adds the URL and runs to three lines here.
  \source{\citeauthor{fisher1926arrangement} (\citeyear{fisher1926arrangement}).
    \citetitle{fisher1926arrangement}.
    \emph{Journal of the Ministry of Agriculture of Great Britain} 33, pp.~503--513.}
\end{frame}

\begin{frame}{Shannon and Tukey: a book opened at random, a slide rule}
  \paperslide{tukey1962future}{\footnotesize
    \begin{itemize}
      \item Shannon's 1948 approximations to English used a book of random
        numbers and two frequency tables; the paper credits Tukey with the
        word ``bit'' \parencite{shannon1948communication}.
      \item In 1956 he warned that the theory had ``been ballooned to an
        importance beyond its actual accomplishments''
        \parencite{shannon1956bandwagon}.
      \item Tukey ``did not use computers'' \parencite{fernholz2003remembering}.
        Screening 36 values by slide rule took two or three times as long as
        a desk calculator and taught two or three times as much; the computer
        was ``important but not vital'' \parencite{tukey1962future}.
      \item Bell Labs' data-analysis group began S on 5 May 1976 to work with
        Tukey's techniques; R replicated S in open source
        \parencite{chambers2020srdatascience}.
    \end{itemize}}
  \source{\fullcite{tukey1962future}}
\end{frame}

\begin{frame}{May and Michener: evidence, and the record-keeping behind it}
  \paperslide{michener2012ecoinformatics}{\footnotesize
    \begin{itemize}
      \item May ``preferred pencil and paper'' \parencite{krebs2020mayobituary};
        his species papers were ``a razor-sharp summary and analysis of the
        available evidence'' \parencite{krebs2022maymemoir}.
      \item In 2011 he set the Library of Congress's count of books
        against our ignorance of species numbers
        \parencite{may2011howmanyspecies}.
      \item The LTER network, from 1980, set aside 15 per cent of each site's
        budget for data management \parencite{stafford2021lterretrospective};
        in 1995 long-term data were still ``only in paper form''
        \parencite{gross1995fled}.
      \item Michener's data life cycle: plan, collect, assure, describe,
        preserve, discover, integrate, analyse. Field data still arrive on
        hand-written sheets \parencite{michener2012ecoinformatics}.
    \end{itemize}}
  % Short, one-line source: a two-line \fullcite runs 3pt into the footer here.
  \source{\citeauthor{michener2012ecoinformatics}
    (\citeyear{michener2012ecoinformatics}).
    \citetitle{michener2012ecoinformatics}.}
\end{frame}

\begin{frame}{An argument from the 1960s}
  \small
  \begin{itemize}
    \item In 1957 Licklider found that about 85 per cent of his ``thinking''
      time went on getting into a position to think
      \parencite{waldrop2002dreammachine}.
    \item His division of labour: people set goals and judge results; the
      machine does the clerical work between decisions.
    \item The information utility of the 1960s was one central machine rented
      by the hour; Unix was small tools joined by pipes. In my view the
      pattern is repeating.
  \end{itemize}
  \fig[0.40\textheight]{ch01-two-lineages}
  \source{Figure from Chapter 1 of the book.}
\end{frame}

\begin{frame}{The model is a function}
  \small
  \begin{itemize}
    \item A language model takes a sequence of \term{tokens} and returns a
      probability for every possible next token. A token is about
      three-quarters of a word; less for scientific names and code.
    \item It has no memory of your last question, no files, no clock and no
      ability to act. Anything that looks like continuity is machinery
      around it.
    \item \term{Pretraining} on a very large body of text; \term{instruction
      tuning} on example answers; \term{preference tuning} towards the answers
      people rank highest, which is where sycophancy is usually traced to;
      \term{reasoning models} add reinforcement learning on tasks a program
      can check.
    \item Hallucination is traced to pretraining, where some false statements
      cannot be told apart from facts, and to evaluations that score a guess
      above ``I don't know'' \parencite{kalai2025hallucinate}.
  \end{itemize}
\end{frame}

\begin{frame}{Chatbot, agent, harness}
  \small
  \begin{itemize}
    \item A chatbot is a model plus a conversation buffer plus a web page:
      good for drafting, explaining and translating; bad at verified fact
      and at acting in the world.
    \item An \term{agent} is a model given tools and a loop. The \term{harness}
      runs the loop: it decides what the model sees, runs the tools, enforces
      permissions and decides when to stop.
    \item Two harnesses running the same model can behave quite differently.
  \end{itemize}
  \fig[0.44\textheight]{ch01-model-chatbot-agent}
  \source{Figure from Chapter 1 of the book.}
\end{frame}

\begin{frame}[fragile]{The whole idea, in pseudocode}
\begin{lstlisting}
while not done and turns < max_turns:   # <- a budget stops it too
    response = model(context)
    context += response                 # the model's own turn
    if response.wants_tool:
        result = run_tool(response.tool, response.args)   # <- permissions here
        context += result
    else:
        done = True
    turns += 1
\end{lstlisting}
  \small
  The \code{run\_tool} line is where a text generator becomes something that
  can delete your files, and where the security risk lives (deck 4).
  \kp{\code{AGENTS.md} is not a wish list. It is a contract that closes
    failures you have seen. Write it after something goes wrong, not before.}
  \source{Listing from Chapter 1 of the book.}
\end{frame}

\begin{frame}{Why now? Three trends met}
  \small
  \begin{itemize}
    \item \textbf{Compute.} Graphics processors, built to draw triangles, are
      fast at exactly what neural networks do: multiplying large matrices in
      parallel. What changed recently is the scale of money. Training a
      frontier model is out of reach for any university; running one on
      hardware you can buy is not (deck 2).
    \item \textbf{Algorithms.} The transformer (next slide), and growth that
      turned out to be predictable: error falls as a power law, the form of
      the species--area curve, with size, data and compute. These
      \term{scaling laws} need ever more text.
    \item \textbf{Data.} Pretraining consumed the accessible text of the
      internet. For ecology the equivalent resource is domain data, curated,
      standardised, with known provenance: the one of the three we own, and
      the subject of Part II.
  \end{itemize}
\end{frame}

\begin{frame}{The transformer, 2017}
  \paperslide{vaswani2017attention}{%
    \begin{itemize}
      \item Introduced the \term{transformer}: \term{attention} combines
        information from different positions in a sequence.
      \item In training, many positions are processed at once, which uses
        GPUs well; the model cannot look ahead at the answer.
      \item Generating still proceeds token by token.
      \item The experiments concerned translation and sentence structure;
        they did not establish ecological expertise
        \parencite{vaswani2017attention}.
    \end{itemize}}
  \source{\fullcite{vaswani2017attention}}
\end{frame}

% ===========================================================================
\section{Powers and limitations (Chapter 3)}

\begin{frame}{What language models are good at}
  \small
  \begin{itemize}
    \item \textbf{Transforming text you supply:} summarising, reformatting,
      translating, prose into structured records. The input constrains the
      output.
    \item \textbf{Extraction:} species interactions
      \parencite{keck2026interactions}, occurrence records in news reports and
      papers. Treat the table as data measured with error: check a sample by
      hand and carry the error rate into the analysis.
    \item \textbf{Code:} the strongest use, because it can be read, tested
      against known answers and rerun. Running is not the same as being right.
    \item \textbf{Getting started and explaining:} a bad first draft beats a
      blank page.
  \end{itemize}
  \medskip
  In common: you supply the facts, or you can check the answer without much
  extra work.
\end{frame}

\begin{frame}{Hallucination: why guessing wins}
  \fig[0.22\textheight]{ch03-guessing-wins}
  \small
  \begin{itemize}
    \item Tests score accuracy, and a model that guesses scores better than
      one that says ``I don't know'' \parencite{kalai2025hallucinate}. We have
      trained good test-takers, and a good test-taker never leaves a blank.
    \item \textcite{zhao2026citations} checked 111 million references: a sharp
      rise in references to papers that do not exist, about 147,000 in 2025
      alone, by their conservative estimate.
    \item Findings age: re-tested on 20 models, the 2023 result that ChatGPT
      invents most references had expired; RAG answers unsupported by the
      evidence persisted \parencite{queen2026shelflife}.
  \end{itemize}
  \kp{Check every reference against a resolver (DOI, PubMed, OpenAlex) before
    it enters a manuscript. Every one, not a sample. A resolver catches
    invented papers, not real ones cited for claims they do not make, so open
    the cited passage too.}
  \source{Figure from Chapter 3 of the book.}
\end{frame}

\begin{frame}{Sycophancy, and the delusional spiral}
  \fig[0.40\textheight]{ch03-sycophancy-loop}
  \small
  \begin{itemize}
    \item Models tend to agree with you: tell them they are wrong and they
      apologise and change their answer, even when your correction is wrong
      \parencite{sharma2024sycophancy}.
    \item Even a perfectly rational user, updating on a source that keeps
      agreeing, can become confidently wrong \parencite{chandra2026spiraling}.
      The input is biased.
    \item It does most damage late in an analysis, alone, with a result you
      want to be true. Each remedy in the figure cuts one arrow.
  \end{itemize}
  \source{Figure from Chapter 3 of the book.}
\end{frame}

\begin{frame}{Five more limits}
  \small
  \begin{itemize}
    \item \textbf{Context.} Everything the model sees in one call is re-sent
      on every turn, and material in the middle is used less reliably. A list
      of taxa fills the window faster than its length suggests. Start a fresh
      session when the task changes.
    \item \textbf{Dependence.} If you no longer write a regression by hand,
      you will not notice when the generated one is wrong. Keep the skill.
    \item \textbf{Provenance.} None: you cannot ask which sources produced a
      claim, or reproduce it exactly. Log prompts, model versions and outputs
      like instrument settings (deck 2).
    \item \textbf{Benchmarks.} Performance is jagged: a system that beats
      people at advanced mathematics can fail an elementary-school problem
      \parencite{mitchell2026jagged}. Test on your own task, with your own
      data, or you do not know (deck 2).
    \item \textbf{Detectors} of AI-written text do not work reliably and are
      biased. Do not use them on people; the answer is disclosure (deck 4).
  \end{itemize}
\end{frame}

% ===========================================================================
\section{Applications in biology and ecology (Chapter 4)}

\begin{frame}{The fullest review for ecologists}
  \paperslide{mammides2026llmecology}{%
    \begin{itemize}
      \item 129 documents on ecology and conservation, 38 of them preprints:
        the evidence is young.
      \item Uses: extraction from text, plain-language database queries,
        literature screening, monitoring news, code, images and sound,
        policy.
      \item Problems: hallucination on less common topics, training-data
        bias, poor reproducibility, paid access, deskilling, energy and
        water.
      \item Recommendations: choose the model for the task, ground with
        retrieval, keep a person checking, publish prompts
        \parencite{mammides2026llmecology}.
    \end{itemize}}
  \source{\fullcite{mammides2026llmecology}}
\end{frame}

\begin{frame}{Extraction at scale: the model is one stage of four}
  \paperslide{keck2026interactions}{%
    \begin{itemize}
      \item 144,402 species interactions from 83,910 articles. Nine in ten
        correct; three in ten real ones missed; the biases of the literature
        inherited \parencite{keck2026interactions}.
      \item Scripts kept paragraphs naming two taxa and an interaction word;
        the model filled a fixed table; names were matched to the NCBI
        Taxonomy; a person checked a sample.
      \item Chapter 11's weekly review has the same shape: scripts on both
        sides of the model (decks 2 and 3).
    \end{itemize}}
  \source{\fullcite{keck2026interactions}}
\end{frame}

\begin{frame}{Two ways to ask a model}
  \fig[0.36\textheight]{ch04-two-ways-to-ask}
  \small
  \begin{itemize}
    \item 100 plant-disease reports: names over 90\% accurate; incidence
      23.8\%, because in 53 reports without a value the model supplied 100\%
      \parencite{gougherty2024extraction}.
    \item 21,955 Red List species: five models placed species in the taxonomy
      correctly 94.9 per cent of the time, and gave the right conservation
      status 27.2 per cent of the time \parencite{uryu2026redlist}.
    \item 200 invented taxon names: models fabricated knowledge for
      0.5--92\% of them \parencite{munch2026phenotypes}. Invented names, which
      no model can have seen, are one control.
  \end{itemize}
  \source{Figure from Chapter 4 of the book.}
\end{frame}

\begin{frame}{The dissent: is ecology ready?}
  \paperslide{mishra2026notready}{\footnotesize
    \begin{itemize}
      \item Ecology is not yet ready for AI, Mishra argues: records contain
        false absences, databases are taxonomically biased, long-term series
        are fragmented \parencite{mishra2026notready}.
      \item Models assume that relationships stay stable, that training data
        represent the future and that observations are independent. Ecology
        violates all three; a model can score well while missing the process.
      \item \textcite{roff2026shape}: readiness is uneven, not absent, and
        underuse has costs. \textcite{mishra2026reply} kept her point:
        detection, classification and monitoring do not necessarily imply
        readiness for inference or prediction under new conditions.
    \end{itemize}}
  \source{\fullcite{mishra2026notready}}
\end{frame}

\begin{frame}{Sampling bias becomes prediction}
  \fig[0.23\textheight]{ch04-sampling-bias}
  \small
  \begin{itemize}
    \item 742 million records: only 6.74 per cent of the globe sampled for
      animals at a 5-kilometre grain; citizen science added to the biases
      \parencite{hughes2021sampling}.
    \item Ecology has methods for uneven sampling: corrections,
      cross-validation for structured data, flags for predictions outside the
      sampled conditions. In my view no language model removes the need for
      them.
    \item BioAnalyst, a foundation model for European biodiversity, is judged
      on an effort-controlled bird survey, not on the GBIF records it learnt
      from, whose biases its authors say it carries
      \parencite{trantas2026bioanalyst}.
  \end{itemize}
  \kp{Before applying a model to an ecological question, ask what would have
    to be true of the data for the answer to be valid, and whether it is.}
  \source{Figure from Chapter 4 of the book.}
\end{frame}

\begin{frame}{The model suggests, the curator decides}
  \paperslide{pombase2026curation}{\footnotesize
    \begin{itemize}
      \item 100 to 200 biocurators maintain the core knowledgebases; the
        literature grows faster than they can read. Their proposal: the
        model suggests, the curator decides \parencite{pombase2026curation}.
      \item GO-Flow: 6,996 articles, 1,826 candidate annotations in 57
        hours; curators agreed with 86.7 per cent of a sample, and review
        took 80 per cent less time.
      \item No AI-extracted association enters a core resource without
        human review. Their ask: an MCP server for every database.
      \item Daisies: 5.8 per cent of filled cells wrong; one cell in six
        differed between runs \parencite{schmidtlebuhn2025morphology}.
    \end{itemize}}
  \source{\fullcite{pombase2026curation}}
\end{frame}

% ===========================================================================
\section{AI-ready data, and protecting your own (Chapters 5 and 6)}

\begin{frame}{What ``AI-ready'' hides}
  \fig[0.21\textheight]{ch05-commons}
  \small
  \begin{itemize}
    \item AI-ready: machine-readable, controlled vocabularies, structured
      metadata, bulk download, open licence. Good data management, done long
      before the term existed.
    \item What the term adds is a consumer that ingests the corpus, builds a
      model, and returns neither the weights nor credit. An argument for
      reading the terms, not against openness.
    \item The Cali Fund (COP16, 2024) asks large companies profiting from
      digital sequence information for 1 per cent of profits or 0.1 per cent
      of revenue, an indicative rate; the sector list ends ``including artificial intelligence''
      \parencite{cbd2024dsidecision}. Sequence data only.
  \end{itemize}
  \kp{When someone asks you to make data ``AI-ready'', ask: ready for whom,
    under what licence, with what attribution, and does anything come back to
    the people who maintain it?}
  \source{Figure from Chapter 5 of the book.}
\end{frame}

\begin{frame}{Open weight is not open source}
  \paperslide{apertus2026report}{\footnotesize
    \begin{itemize}
      \item \textbf{Open weight:} download, run, fine-tune. Not: see the
        training data, or rebuild it.
      \item \textbf{Open source} (OSI): weights, the full training and
        inference code, and enough about the data to build an equivalent
        \parencite{osi2024definition}.
      \item OLMo released data, code and logs \parencite{groeneveld2024olmo}.
        Apertus, from Swiss public institutions: only openly available data,
        crawler opt-outs honoured retroactively, Apache 2.0
        \parencite{apertus2026report}.
      \item Without the training data nobody can audit which regions, taxa
        and languages are under-represented
        \parencite{hai2026openweight,torresani2026opening}.
    \end{itemize}}
  \source{\fullcite{apertus2026report}}
\end{frame}

\begin{frame}{What makes biodiversity records fit for a model}
  \paperslide{wieczorek2012darwincore}{\footnotesize
    The format can be right and the inference wrong. Six kinds of information
    decide what a record supports \parencite{wieczorek2012darwincore}:
    \begin{itemize}
      \item the sampling event and its effort (\code{samplingEffort});
      \item non-detections (\code{occurrenceStatus} \code{absent});
      \item coordinate uncertainty in metres;
      \item the taxonomic backbone and its version;
      \item identification qualifiers (``cf.'', \code{identifiedBy});
      \item flags for machine-derived records, unverified until a person
        checks.
    \end{itemize}
    The springtail survey of Chapter 10 keeps a non-detection and an effort
    column for this reason. A model cannot ask the collector what was left
    out.}
  \source{\fullcite{wieczorek2012darwincore}}
\end{frame}

\begin{frame}{Standards, and a protocol for machines}
  \fig[0.36\textheight]{ch05-mcp}
  \small
  \begin{itemize}
    \item Darwin Core carries occurrence exchange; MIxS asks for a sample's
      origin and processing alongside the sequence
      \parencite{wieczorek2012darwincore}; ENVO and EUNIS do the semantic work.
    \item FAIR data still needs harmonising before it can feed a model: units,
      methods, current ontology versions. Today mostly manual and seldom
      reproducible \parencite{woodcharlson2026cope}.
    \item MCP wraps a resource once instead of once per client. A server runs
      with your access, and its tool descriptions are text the model may obey
      (deck 4).
  \end{itemize}
  \source{Figure from Chapter 5 of the book.}
\end{frame}

\begin{frame}{Retrieval-augmented generation, and its seven failure points}
  \fig[0.46\textheight]{ch05-rag-pipeline}
  \small
  \begin{itemize}
    \item Ask ``what colour is my ball?'' and a model invents an answer; say
      ``I have a green box with a red ball in it'' first and it answers
      correctly. RAG automates that second clause.
    \item Three failures are in retrieval and four in generation
      \parencite{kurniawan2024ragfailures}. Log what was retrieved separately
      from what was generated; test retrieval by itself first; for a small,
      stable corpus, notes in the context can beat retrieval.
  \end{itemize}
  \source{Figure from Chapter 5 of the book.}
\end{frame}

\begin{frame}[fragile]{A research tree, PARA-shaped}
\begin{lstlisting}[style=shell]
~/work/
  10_projects/2026_posidonia_edna/
    README.md   AGENTS.md
    data/raw/             # never edited, read-only
    data/raw.sha256       # checksums of data/raw, written before it is locked
    data/samples.tsv      # the sample sheet: one row per sample, the record
    data/accessions.tsv   # archive accessions of the reads
    data/processed/       # regenerable from raw + scripts
    scripts/  results/  manuscript/
  20_areas/monitoring_network/
  30_resources/literature/
  40_archive/2024_thermal_tolerance/
\end{lstlisting}
  \small
  \begin{itemize}
    \item An agent finds its way by reading the files; \code{final2/} and
      \code{Untitled1.csv} leave it guessing.
    \item Raw reads belong in a public archive; the project keeps what points
      at them. The sample sheet is the record.
  \end{itemize}
  \source{Listing from Chapter 6 of the book.}
\end{frame}

\begin{frame}[fragile]{Raw data is read-only}
\begin{lstlisting}[style=shell]
find data/raw -type f -print0 | sort -z | xargs -0 shasum -a 256 > data/raw.sha256
chmod -R a-w data/raw/
shasum -a 256 -c data/raw.sha256     # verify later: every line OK, exit status 0
\end{lstlisting}
  \fig[0.23\textheight]{ch06-raw-readonly}
  \small
  \begin{itemize}
    \item Never modify raw data: checksums, a read-only folder, everything
      else derived by script. Lock first, or the checksums cannot be written;
      keep \code{raw.sha256} beside the folder.
    \item Agents can delete files \parencite{shapira2026agentsofchaos}.
      Read-only permissions turn in-place ``cleaning'' into a visible error:
      a tripwire, not a wall; the checksums say afterwards what changed.
    \item The workshop's rules say it to the agent: ``\code{data/raw/} is
      READ-ONLY. Never edit, move or overwrite files there.''
  \end{itemize}
  \source{Listing and figure from Chapter 6 of the book.}
\end{frame}

\begin{frame}{Backup 3-2-1, Git, and plain text}
  \fig[0.38\textheight]{ch06-backup}
  \small
  \begin{itemize}
    \item Three copies, two kinds of storage, one off-site; restore a random
      file from each copy every quarter. RAID is not a backup: deletion,
      ransomware and \code{rm -rf} reach every disk at once.
    \item Git for code, always; with agents it stops being optional. Commit
      before an agent starts, push the checkpoint, read \code{git diff}
      afterwards; deny \code{git reset --hard} and \code{git push --force}
      in the harness.
    \item Plain text lasts, and agents handle it best; sequence data keep
      their own formats (FASTQ, BAM, CRAM, HDF5).
  \end{itemize}
  \source{Figure from Chapter 6 of the book.}
\end{frame}

% ===========================================================================
\begin{frame}{Takeaways}
  \small
  \begin{itemize}
    \item The machine takes the drudgery, and the judgement stays with the
      ecologist.
    \item Names will change; the design (model, loop, tools, permissions,
      files) has been stable, and is the part worth learning.
    \item We have trained good test-takers, and a good test-taker never
      leaves a blank. Check every reference against a resolver: every one,
      not a sample.
    \item The question is not whether the method is powerful but whether the
      data can support the inference.
    \item Ecology's data commons was built by people, over decades, to be
      shared. ``AI-ready'' describes what that work made possible, not who did
      it or on what terms it is now being consumed.
    \item Never modify raw data. Record checksums, make the folder read-only,
      and derive everything else by script.
  \end{itemize}
\end{frame}

\begin{frame}[allowframebreaks]{References}
  \printbibliography[heading=none]
\end{frame}

\end{document}
